\documentclass[a4paper,11pt]{article}
\usepackage{exam-style-842-unit}
\examinehead{842 信息化概论 · 第06单元练习册}{练习三 参考答案与解析}
\begin{document}

\unitanswerscover{842 信息化概论 · 单元练习册}{第06单元参考答案}{练习三（物联网应用与新一代信息技术）的答案与解析；每道判断题均给出改正后的正确表述；含安全与隐私、标准体系的答题要点。}

\answerpart{一、单项选择题（10 题，每题 3 分，共 30 分）}

\q A。\begin{analysis}教材 7.1.3 指出，对牲畜溯源是给放养的牲畜中的每一只都贴上一个 RFID 电子标签（电子耳码），这个电子耳码会一直保留在超市出售的肉品上，消费者可通过 RFID 阅读器阅读电子耳码，知道牲畜的成长历史，确保食品安全（教材 7.1.3，第 116 页）。B 二维码也可用于溯源，但牲畜佩戴的 RFID 标签教材称为电子耳码；C 电子鼻是气体传感器；D 电子秤是称重设备。\end{analysis}

\q A。\begin{analysis}教材 7.1.3 指出，最典型的精准农业应用是 2002 年英特尔公司率先在俄勒冈建立的世界上第一个无线葡萄园：将传感器节点分布在葡萄园的每个角落，每隔一分钟检测一次土壤温度、湿度或该区域有害物的数量，以确保葡萄健康生长（教材 7.1.3，第 116 页）。B、C、D 均不是建立该无线葡萄园的公司。\end{analysis}

\q B。\begin{analysis}教材 7.1.3 指出，上海浦东国际机场防入侵系统铺设了 3 万多个传感节点，覆盖了地面、栅栏和低空探测，多种传感手段组成一个协同系统，可以防止人员的翻越、偷渡、恐怖袭击等攻击性入侵（教材 7.1.3，第 116 页）。A、C、D 的数量均与教材不符。\end{analysis}

\q B。\begin{analysis}停车场案例中，出入口地磁传感器感知到车辆出入行为时，启动摄像机拍照，通过车牌识别器识别车辆牌照信息，并将牌照信息、出入行为、出入时间和道闸编号等信息通过移动互联网络传输至中心管理平台（教材 7.1.4，第 117 页）。A 道闸由管理平台的控制指令驱动；C 信号灯由车位超声传感器控制；D 电视大屏位于监控中心。\end{analysis}

\q A。\begin{analysis}教材 7.1.4 指出，各车位上的超声传感器采用 ZigBee 短距离无线通信组网，形成覆盖整个停车场的传感器网络，实时将空闲车位信息通过通信终端传输至中心管理平台（教材 7.1.4，第 117 页）。B 卫星通信用于广域覆盖；C 光纤通信是有线方式；D 广播电视网是单向广播网络，均不是车位传感器组网技术。\end{analysis}

\q B。\begin{analysis}教材 7.2 明确指出：三网融合主要指电信网、移动互联网以及广播电视网的融合，此融合并非三网的物理融合，而是应用上的有机融合（教材 7.2，第 121 页）。A 把融合理解为物理网络合并，与教材相反；C、D 均不是教材对三网融合的界定。\end{analysis}

\q D。\begin{analysis}教材 7.2 指出，新一代信息技术包括 6 个方面：下一代通信网络、物联网、三网融合、新型平板显示、高性能集成电路和云计算（教材 7.2，第 121 页）。A 物联网、B 三网融合、C 云计算均在其中；D 传统机械制造不属于新一代信息技术。\end{analysis}

\q D。\begin{analysis}教材 7.1.5 指出，物联网的大规模应用仍面临巨大挑战，至少包括成本的挑战、安全的挑战、侵犯隐私的威胁 3 个方面（教材 7.1.5，第 120 页）。D 错误：教材同时指出目前物联网还缺乏统一标准、标准比较分散、体系还不完善，“标准完全统一”是尚未实现的目标而非已列出的挑战，且与“标准分散”的现状相反。\end{analysis}

\q A。\begin{analysis}教材第 11 章指出，就总体目标而言，家庭自动化就是要给人们创造一个安全、高效、优雅、舒适、节能和环保的居住环境；就技术手段而言，要充分利用计算机、通信、检测、人工智能、自动控制、机器人等现代科技管理和控制家中各种设备（教材第 11 章，第 207 页起）。B、C、D 均不是教材表述。\end{analysis}

\q A。\begin{analysis}教材 7.1.5 指出，欧洲智能系统集成技术平台组织（EPoSS）在《Internet of Things in 2020》中预测，物联网的发展将经历 4 个阶段：2010 年之前以 RFID 为代表的技术广泛应用于物流、零售和制药等领域；2010—2015 年实现物与物之间的互联；2015—2020 年进入半智能化；2020 年之后实现全智能化（教材 7.1.5，第 120 页）。B 美国 FORRESTER 公司预测的是 2020 年物物互联业务与人际通信业务之比将达 30:1；C ITU 报告提出的是物联网概念拓展到传感网络；D 中投顾问报告给出的是我国物联网市场规模。\end{analysis}

\answerpart{二、多项选择题（4 题，每题 4 分，共 16 分）}

\q ABC。\begin{analysis}教材 7.1.3 所列的物联网典型应用有 4 个：牲畜溯源、精准农业、安全防入侵系统、医疗保健（个人保健）（教材 7.1.3，第 116 页）。A、B、C 正确。D 错误：数控机床加工属于制造装备领域，不是教材 7.1.3 所列的物联网典型应用。\end{analysis}

\q ABC。\begin{analysis}教材 7.1.5 把我国物联网发展现状概括为 5 个方面：产业体系初步形成，但产业化能力不高，尚未形成规模化产业优势；核心关键技术有待突破，在传感器、芯片、关键设备制造、智能通信与控制、海量数据处理等核心技术上与发达国家存在较大差距；标准比较分散，体系还不完善，在国际上面临核心标准的竞争；物联网应用的规模和领域比较小，没有形成成熟的商业模式，应用成本较高；承载大量国家经济社会活动和战略性资源，面临巨大的安全与隐私保护挑战（教材 7.1.5，第 120 页）。A、B、C 正确。D 错误：与教材“应用规模和领域比较小、没有形成成熟商业模式”正好相反。\end{analysis}

\q ABC。\begin{analysis}教材 7.2 指出，新一代信息技术包括下一代通信网络、物联网、三网融合、新型平板显示、高性能集成电路和云计算 6 个方面（教材 7.2，第 121 页）。A 下一代通信网络、B 物联网、C 三网融合均在列。D 错误：传统机械制造不属于新一代信息技术。\end{analysis}

\q BCD。\begin{analysis}教材 7.1.5 指出，物联网大规模应用至少面临 3 个方面的挑战：成本的挑战、安全的挑战、侵犯隐私的威胁（教材 7.1.5，第 120 页）。B、C、D 正确。A 错误：“标准完全统一”不是教材列出的挑战，教材强调的恰恰是目前还缺乏统一标准、标准比较分散。\end{analysis}

\answerpart{三、判断题（5 题，每题 2 分，共 10 分）}

\q √。\begin{analysis}教材 7.1.3 对精准农业的定义是：研究基于农田、园林、温室等目标区域的射频技术与无线传感器网络架构，实现大量实时地收集温度、湿度、光照、气体浓度等物理量，精准地获取土壤水分、压实程度、电导率、pH 值、氮素等土壤信息，实现精细化的农业生产控制（教材 7.1.3，第 116 页）。\end{analysis}

\q √。\begin{analysis}教材 7.1.2 指出，应用层是物联网系统结构的最高层，包括各类用户界面显示设备及其他管理设备；应用层根据用户需求可面向各类行业建立实际应用的管理平台和运行平台，如绿色农业、公共安全、远程医疗、智能交通等（教材 7.1.2，第 115 页）。\end{analysis}

\q ×。改正：目前物联网还缺乏统一标准，我国标准比较分散、体系还不完善，国际上面临核心标准的竞争；我国虽已成立电子标签 RFID 标准工作组、传感器网络标准工作组、泛在网技术工作委员会和中国物联网标准联合工作组，但统一的标准体系尚未建成。\begin{analysis}教材 7.1.5 在“物联网的标准化概况”中明确指出：标准化的实现能够整合行业应用、保证互操作性和互联互通，是扩大市场规模的基础；但目前物联网还缺乏统一标准（教材 7.1.5，第 120 页）。把“已完善、已统一”当作事实，是本单元标准部分最常见的错误。\end{analysis}

\q ×。改正：三网融合并非三网的物理融合，而是电信网、移动互联网以及广播电视网在应用上的有机融合。\begin{analysis}教材 7.2 原文为：“此融合并非三网的物联（物理）融合，而是应用上的有机融合”（教材 7.2，第 121 页）。混淆“物理合一”与“应用融合”是三网融合的经典辨析点。\end{analysis}

\q ×。改正：智慧农业是以物联网、云计算、大数据、人工智能等新一代信息技术为支撑，实现农业生产全过程“感知—决策—控制”闭环和全链条管理的现代农业形态；安装摄像头远程监控只是其中一个环节，不能等同于智慧农业。\begin{analysis}智慧农业强调用信息技术对农业生产、经营、管理、服务全过程进行感知、决策与控制，摄像头监控只解决“看得见”，还需要传感器采集、数据分析和智能决策（联合教材 7.1.2 应用层“绿色农业”与 7.1.3 精准农业的理解作答；教材未单列“智慧农业”定义，按大纲与通行表述补充）。\end{analysis}

\answerpart{四、配对题（4 对，每对 1 分，共 4 分）}

\q （1）—B\quad （2）—A\quad （3）—D\quad （4）—C。
\begin{analysis}精准农业的典型是 2002 年英特尔建立的无线葡萄园；牲畜溯源使用 RFID 电子耳码；安全防入侵的典型是上海浦东国际机场 3 万多个传感节点；停车场中心管理平台由通信服务器、管理服务器、计算服务器、数据服务器、媒体服务器和应用服务器等组成。（教材 7.1.3、7.1.4，第 116—117 页）\end{analysis}

\answerpart{五、简答题（4 题，每题 6 分，共 24 分）}

\q 参考答案要点：
（1）牲畜溯源（3 分）：给放养的每头牲畜贴一个 RFID 电子标签（电子耳码），电子耳码记录牲畜身份与成长历史，并一直保留到超市出售的肉品上；消费者或管理者用 RFID 阅读器读取耳码，即可追溯牲畜的养殖、防疫、屠宰、流通信息，实现食品安全追踪溯源。目前新疆、宁夏、江苏等地已广泛应用。技术实现上依靠 RFID 非接触识别、电子标签存储与后端信息系统。（教材 7.1.3，第 116 页）
（2）精准农业（3 分）：基于农田、园林、温室等目标区域的射频技术与无线传感器网络架构，实时收集温度、湿度、光照、气体浓度等物理量，精准获取土壤水分、压实程度、电导率、pH 值、氮素等土壤信息，实现精细化的农业生产控制。最典型的案例是 2002 年英特尔在俄勒冈建立的世界上第一个无线葡萄园：传感器节点分布到每个角落，每分钟检测土壤温度、湿度或有害物数量；通过长年数据记录与分析，掌握葡萄品质与日照、温度、湿度的确切关系，指导葡萄生长与葡萄酒品质控制。目前物联网在精细农业的大规模应用还处于起步阶段，示范工程已有一定应用。（教材 7.1.3，第 116 页）
（3）常见失分：只写“溯源”“精准”概念而不给技术手段；把无线葡萄园的公司、年份写错。

\q 参考答案要点（教材 7.1.4，第 117 页）：
（1）感知设备（2 分）：出入口地磁传感器感知车辆出入行为，启动摄像机拍照并由车牌识别器识别牌照；车位超声传感器探测车辆进入或离开车位，控制车位上方信号灯亮红灯或绿灯，并采用 ZigBee 短距离无线通信组网，把空闲车位信息实时上传。感知设备完成车辆、车位、时间等信息的采集。
（2）传输网络（1 分）：随着 3G、4G 移动通信技术的发展，移动通信网络把分布在各地区的停车场感知设备互联至互联网，形成物联网网络；感知设备终端与互联网平台中的通信服务器通常采用 TCP/IP 互联并交换数据。
（3）中心管理平台（2 分）：通常部署在互联网云计算平台或停车场业主计算机机房内，由通信服务器、管理服务器、计算服务器、数据服务器、媒体服务器和应用服务器等组成；负责数据存储、计算处理、业务管理与服务提供，经计算处理后提示车辆计价及通行控制。
（4）监控中心（1 分）：通常由测控计算机和电视大屏组成，为停车场的运营管理和安全保障提供运营和管控决策信息。
（5）答题提示：四部分对应“感知—传输—支撑/平台—应用”的物联网逻辑，可联系教材四层架构说明：感知设备属感知层，传输网络属网络层，中心管理平台属支撑层，监控中心属应用层。

\q 参考答案要点：
（1）新一代信息技术 6 个方面（2 分）：下一代通信网络（NGN）、物联网、三网融合、新型平板显示、高性能集成电路、云计算。其中 NGN 建立在高带宽光纤网络和 3G 无线网络基础上，实现音频、视频、数据统一；三网融合是电信网、移动互联网、广播电视网在应用上的有机融合；新型平板显示以 OLED 和电子纸为代表；云计算把计算任务分布到资源池上，按需提供计算、存储和软件服务。（教材 7.2，第 121 页）
（2）物联网与云计算（1.5 分）：物联网负责感知与连接，产生海量数据；云计算提供弹性、可扩展的计算与存储资源，是物联网数据的存储与处理平台，是物联网发展重要的技术支撑。教材 7.1.2 把云计算列为支撑技术，并指出云计算包含基础设施即服务（IaaS）、平台即服务（PaaS）和软件即服务（SaaS）三个层次。（教材 7.1.2，第 115 页）
（3）物联网与大数据（1.5 分）：物联网的传感器、RFID、视频等持续产生海量、多源、异构数据，为大数据分析提供数据源；大数据技术负责对这些数据进行存储、清洗、分析与挖掘，从数据中发现规律，二者是“数据源”与“数据加工”的关系。（教材 7.1.2 把数据库与数据挖掘列为支撑技术；大数据概念教材第 4 章展开）
（4）物联网与人工智能（1 分）：物联网提供实时、多维的感知数据，人工智能负责识别、预测与决策，把数据转化为智能控制指令，实现病虫害识别、灌溉决策等；教材 7.1.2 把人工智能技术列为物联网的支撑技术之一。（教材 7.1.2，第 115 页）
（5）结论：物联网、云计算、大数据、人工智能不是彼此替代，而是“感知—传输—存储计算—分析挖掘—智能决策”的协同链条，共同构成新一代信息技术体系。

\q 参考答案要点：
（1）发展现状（2 分）：教材概括为 5 个方面：产业体系初步形成，但产业化能力不高，尚未形成规模化产业优势；核心关键技术有待突破，在传感器、芯片、关键设备制造、智能通信与控制、海量数据处理等核心技术上与发达国家存在较大差距；标准比较分散，体系还不完善，在国际上面临核心标准的竞争；物联网应用的规模和领域比较小，没有形成成熟的商业模式，应用成本较高；物联网承载大量国家经济社会活动和战略性资源，面临巨大的安全与隐私保护挑战。（教材 7.1.5，第 120 页）
（2）标准化概况（1.5 分）：国际方面，ETSI 欧洲电信标准化协会的 M2M 技术委员会研究 M2M 物联网，ITU-T 研究泛在传感网（USN）标准，ISO/IEC 研究 RFID 标准化；我国 2005 年 12 月成立电子标签 RFID 标准工作组，2009 年 9 月成立传感器网络标准工作组（WGSN），2010 年 2 月中国通信标准化委员会成立泛在网技术工作委员会，2010 年 6 月成立中国物联网标准联合工作组（下设 19 个专业委员会）。目前物联网还缺乏统一标准，标准体系建设是扩大市场规模的基础和产业发展的关键。（教材 7.1.5，第 120 页）
（3）产业前景（1.5 分）：EPoSS 在《Internet of Things in 2020》中预测物联网将经历 4 个阶段，2020 年之后实现全智能化，网络从虚拟走向现实、从局域走向泛在，在 anytime、anyone、anywhere 基础上拓展到 anything；美国 FORRESTER 预测到 2020 年物物互联业务与人际通信业务之比将达 30:1，物联网被称为下一个万亿美元级的通信业务；我国物联网产业将出现井喷式发展，释放上万亿元市场容量。（教材 7.1.5，第 120—121 页）
（4）主要挑战（1 分）：成本的挑战、安全的挑战、侵犯隐私的威胁。（教材 7.1.5，第 120 页）
（5）常见失分：把“标准已统一”当作现状；只答前景不答挑战；混淆 EPoSS 与 FORRESTER 的预测结论。

\answerpart{六、综合题（1 题，共 16 分）}

\q 参考答案要点：
（1）物联网在各环节的作用（6 分，每环节 1.5 分）：
种植环节——气象站与土壤传感器实时采集温度、湿度、光照、土壤水分、pH 值、氮素等数据，北斗终端提供农机位置与作业轨迹，实现精准灌溉、精准施肥和农机调度；
加工环节——在稻米加工包装环节赋二维码和 RFID 标签，记录加工批次、加工时间、质检结果，实现批次管理与质量追溯；
物流环节——RFID/二维码在仓储、运输节点被读写器自动采集，结合北斗定位实现运输路径与温湿度监控，防止串货和变质；
销售环节——消费者扫码查询产地、农事记录、检测报告，实现“来源可查、去向可追、责任可究”，提升品牌信任与产品溢价。
上述环节体现了物联网“全面感知—可靠传输—智能处理—行业应用”的闭环，也对应教材应用层的绿色农业、智能交通（物流）等应用平台。
（2）云计算、大数据、人工智能的角色与协同（6 分）：
云计算（2 分）——承担存储与计算基础设施角色，把海量传感器数据、影像、溯源记录集中存储，按需提供计算资源，是平台的“底座”；对应教材支撑层的云计算设备与超级平台。
大数据（2 分）——承担数据分析角色，对多源、异构、长时序的产量、气象、土壤、市场数据进行分析挖掘，找出产量影响因素、优化种植方案；对应教材支撑技术中的数据库与数据挖掘技术。
人工智能（2 分）——承担智能决策角色，用模型识别病虫害、预测发生趋势并推送防治方案，把数据转化为可执行的农事指令；对应教材支撑技术中的人工智能技术。
协同关系：物联网负责“采得到、传得回”，云计算负责“存得下、算得动”，大数据负责“挖得深”，人工智能负责“判得准”，四者构成“感知—传输—存储计算—分析挖掘—智能决策”的完整链条，缺一不可。
（3）改进建议（4 分，答出 3 条即可，每条要落到具体措施）：
安全与隐私方面——建立数据分级分类与访问控制制度，对农户身份、地块位置、企业经营数据等敏感信息脱敏或加密存储；溯源数据上链或采用数字签名防篡改，明确数据采集、使用、共享的授权边界；落实网络安全等级保护，定期开展漏洞扫描与应急演练。
标准体系方面——对接国家物联网标准体系，统一传感器接口、编码标识（如 EPC/二维码编码规则）、数据格式与平台接口，避免“一县一标准、一企一平台”；参照国际 ETSI M2M、ITU-T USN、ISO/IEC RFID 标准与我国传感器网络标准工作组（WGSN）成果，制定地方标准并与行业标准、国家标准衔接；建立标准符合性测试与互认机制，降低多平台对接成本。
产业与推广方面——先做示范园区再规模化推广，探索“政府引导＋企业运营＋农户付费”的商业模式，降低应用成本；培育本地运维与数据服务队伍，解决“建得起、用不好”的问题。
（4）结论（加分点）：该平台是把物联网、云计算、大数据、人工智能集成的智慧农业与溯源系统，体现了新一代信息技术“感知—连接—计算—智能”的协同；答题时按“环节作用—技术分工—安全标准建议”三层展开，可覆盖 2021 年分析题“物联网技术如何用于农产品溯源”、2023 年论述题“农业物联网的关键技术及在智慧农业中的应用”的采分点。（教材 7.1.2、7.1.3、7.1.5、7.2，第 114—121 页）

\end{document}
